\section{Durability Across a Remote API}
\label{sec:remote}

The harness above was validated on local GPU jobs, where the failure modes are the
process, the disk and the write syscall. A remote model adds failure classes that
none of those tests reach: the provider can return a well-formed HTTP 200 that
carries no answer, close TLS connections under concurrency, or refuse service for
billing reasons while the client is healthy. We ran the same durable contract
against a closed-weight model through an OpenAI-compatible relay
(\texttt{deepseek-v4.1-flash}) for 9,400 calls, using the same prompts, parser,
single-writer queue and provenance fingerprint as the local arm.

\begin{table}[t]
\centering
\footnotesize
\setlength{\tabcolsep}{3.5pt}
\caption{Failure classes observed only on the remote arm. ``Rows'' is the number
of cells each class touched in the delivered 9,400-call run. Every class was
detected by the harness and none produced a silently wrong row.}
\label{tab:remote}
\begin{tabular}{p{0.20\textwidth}p{0.20\textwidth}p{0.19\textwidth}p{0.19\textwidth}r}
\toprule
Class & Observed & Detection & Recovery & Rows \\
\midrule
Hidden reasoning exhausts the budget & HTTP 200, \texttt{choices: []},
\texttt{completion\_tokens}\allowbreak${=}$cap & empty choice list kept as
\texttt{finish\_reason}\allowbreak\texttt{=}\allowbreak\texttt{empty\_choices},
not a transport error & raise the cap to
8{,}192; record residual cells as contract failures & 5 \\
TLS handshake reset under concurrency & \texttt{EOF occurred in violation of
protocol}; throughput fell to 1 row/min & transport exception, retried with
backoff & one keep-alive connection per worker thread; 11--25 rows/min & --- \\
Provider quota wall & HTTP 402
\texttt{insufficient\_}\allowbreak\texttt{balance} after 7,315
successful cells & status code recorded per row & \texttt{retry\_failed.py}
removes failed rows and the sidecar; rerun after reset & 2{,}085 \\
Provenance drift across restarts & 149 cells written by an earlier runner
revision & per-row provenance id & keep the majority id, rerun the minority &
149 \\
\bottomrule
\end{tabular}
\end{table}

\paragraph{The invariants extend, but not unchanged.} The local contract's fifth
invariant (``no silent configuration mixture'') already covers the provenance
drift. Two further statements are needed for a remote provider.

\begin{enumerate}[leftmargin=1.4em,itemsep=1pt,topsep=2pt]
\item \textbf{HTTP success is not answer success.} A 200 response with an empty
choice list is a contract failure and must be recorded as one. Without this rule
the arm would have reported 1,024-token cells as ``answered with no findings'',
inflating every denominator.
\item \textbf{Recovery is provenance-scoped.} A resumed stage may only absorb rows
written under the current fingerprint. The 149 stale cells were not ``reused with
a note''; they were rerun, so the delivered file carries exactly one provenance
id and 9,400 unique keys.
\end{enumerate}

\paragraph{What the remote run cost.} 9,400 calls, 0 transport errors in the
delivered file, and one quota interruption. The interruption is the interesting
one for a benchmark paper: it is not a fault in the scientific design, but it is
the most likely way a large evaluation silently loses cells. The recovery
procedure is mechanical and auditable --- the failed rows and their byte offsets
are removed, the file is backed up with a timestamp, and the rerun refills exactly
those cells --- which is what makes the loss reportable rather than invisible.

\paragraph{Limits.} One provider and one model. The relay also hides its reasoning
channel, so the output-mode invariant of \S\ref{sec:evaluation} cannot be checked
remotely; we report the contract axis only.
